\section{Discussion}
\label{sec:discussion}
This section discusses the applicability of \thename's unikernel-based oblivious paging scheme on other \gls{cvm} architectures like Intel TDX and ARM CCA and provides an analysis of \thename security guarantees.


% \subsection{Applicability to other \gls{cvm}}
\label{sec:discussion:cvms}

% \hdr{Prtoection mechanisms }
% While Intel TDX's implementation is different, the core principle of guest VM protection is akin to SEV-SNP.
% TDX platforms provide nested paging-based memory management, allowing direct MMU access to guests.
% TDX provides transparent encryption and guarantees the integrity of GPA to HPA mapping with PAMT, similar to SEV-SNP's RMP table.
% Also, TDX requires guest TD to accept(\texttt{TDG.MEM.PAGE.ACCEPT}) and track pages.

\iffalse
\hdr{Controlled-channel threats on Intel TDX and ARM CCA}
Our analysis of the publically available specifications for Intel \gls{tdx}~\cite{tdx,intel-tdx-survey} and ARM \gls{cca}~\cite{cca} indicates that while ARM \gls{cca} is free from controlled-channels attacks, Intel \gls{tdx} remains vulnerable.
ARM \gls{cca} currently supports only restricted dynamic memory management: once a \gls{cvm} page is unmapped, \gls{cca} discards its content and mapping, preventing re-mapping~\cite{cca}.
This approach mitigates \gls{npf} controlled-channel attacks but results in less flexible memory management.
Intel \gls{tdx} prevents direct \gls{npt} manipulation, but the hypervisor could still intercept memory accesses through the \emph{page blocking} mechanism~\cite{tdx}.
\fi


\hdr{Applicatibity to other architectures}
\thename's security policies have been developed for AMD SEV-SNP.
However, they can be generalized to other TEE architectures.
For instance, we expect SGX enclaves to also benefit from our work, but we leave it as future work.
In the case of Intel \gls{tdx},  we found that in its current specification, even without direct \gls{npt} access, the hypervisor could still perform \gls{npf} controlled-channel attacks through the \emph{page blocking} mechanism~\cite{tdx}.
In addition, \gls{tdx} provides a minimum of 1,000 CVM execution cycles between each hypervisor-induced VMExit, thereby reducing the temporal resolution of single-stepping attacks.
Also, it adds random variance to the minimum cycles each time to add noise to the side-channel attack launched on top of single-stepping.
Thus, \thename may offer the same level of security as its SEV implementation with a relaxed tick rate.


% However, such a management scheme severely limits the hypervisor's ability to manage memory. 
% Thus, in the current
% CCA does not support deallocating preallocated guest VM pages, lacking flexible memory management.
% It only supports lazy allocation of guest pages on the fly.
% The host can destroy the pages and reclaim them; however, it cannot repopulate them with previous data; only the terrorized page can do this.
% 
% Every allocated guest page is managed within \gls{rmm}, revealing no trace of page faults to the host.
% However, CCA, mostly being enforced with firmware, can change to support these features in the future.


% \hdr{Avaiability of \thename's building blocks on other CVMs}
% \thename leverages huge mappings and direct memory access
% % Also, both platforms allow huge page mapping up to 1GB pages~\cite{}.
% Both \gls{tdx} and \gls{cca} support huge page mappings and provide allowing direct MMU access to \glsps{cvm}.
% Moreover, both platforms provide a secure timer source to implement \thename's validation scheme.
% % (But CCA doesn't support it currently)
% We expect that \thename can be adapted to any \gls{cvm} platform without major design changes.

\hdr{Trusted timer interrupt}
Our binary instrumentation to achieving self-sovereignty may be replaced with preemptive scheduling if a trusted source of timer interrupts exists.
AMD Secure AVIC specification \cite{secureavic} secure guest's \gls{apic} states and prevents the hypervisor from injecting unexpected interrupts into \glspl{cvm}.
Unfortunately, while this shows potential in securing interruptions, the current Secure AVIC specification still requires the hardware timer (LAPIC) to be emulated by the untrusted hypervisor~\cite{secureavic}.
% Thus, it is a future research direction to enable a trusted source of timer that respects deadline constraints for \gls{cvm}.

\hdr{\textcolor{red}{Kill discussion}}


% **Note TDX have single stepping protection + continous faulting GPA protection for NPF attacks**

% \hdr{Security of bounded leakage}
% While we employ the strictest bounded leakage $n=1$ in our evaluation, in practice, this value might be relaxed to reduces the performance overheads.

% This is because on \gls{sevsnp}, once a page is smashed, its merging requires 

% \subsection{Limitations}



% \hdr{Binary rewriting}
% \thename inherits limitations of binary rewriting tools.
% \thename employs probe-based binary rewrite to avoid complications with control flow recovery and relocation, two major issues with binary rewriting~\cite{e9patch}.
% However, most rewriting tools also depend on correct program disassembly, which may be challenging in certain cases, for instance, data embedded within code.
% Thus, \thename may
